\section{Introduction}
\label{sec:intro}

Whether the aligned behavior of a large language model (LLM) is encoded
\emph{redundantly} decides whether safety survives quantization, pruning, benign
fine-tuning, or adversarial attack. This is not a matter of taste: the same
intervention that is harmless to a redundantly encoded behavior can be fatal to one
that depends on a single internal component. Yet the empirical literature on the
internal mechanisms of alignment is sharply, and at first sight irreconcilably,
divided.

On one side, alignment looks \emph{fragile and localized}. Fine-tuning on fewer than
a hundred examples raises harmful-response rates from near zero to roughly
$80$--$90\%$ in five to ten gradient steps \cite{qi2023finetuning}; the divergence
between an aligned model and its base model is concentrated in the first few output
tokens \cite{qi2024shallow}; and a single difference-in-means direction, when
projected out, removes refusal across more than a dozen models \cite{arditi2024refusal}.
On the other side, alignment looks \emph{distributed and robust}. Ablating all three
primary name-mover heads of the indirect-object-identification circuit drops the
relevant logit difference by only $5\%$, because dormant backup heads activate to
compensate \cite{wang2022ioi}; refusal is mediated not by one direction but by
polyhedral concept cones of dimension up to five \cite{wollschlager2025cones}; and a
model can survive ablation of its principal safety direction because of ``a richer
internal refusal structure with redundant paths that survived the ablation''
\cite{joad2026directions}.

\paragraph{The gap.}
What is missing is a single quantitative account under which both regimes are special
cases---a scalar (or a pair of scalars) that says \emph{which} kind of robustness a
given aligned behavior has, and which it lacks. The canonical formal measure of
compensation, McGrath et al.'s logit-level compensatory effect \cite{mcgrath2023hydra},
is layer-granular and, as Rushing and Nanda show, roughly $30\%$ of it is a
LayerNorm-scaling artifact rather than genuine rerouting \cite{rushing2024selfrepair}.
No assumption-light bound currently links a redundancy measure to the
\emph{behavioral} effect of an ablation, and no model places ``shallow/fragile'' and
``distributed/redundant'' on one axis.

\paragraph{Our contribution.}
We supply such a framework, organized in three proved layers, each verified
computationally.

\begin{itemize}[leftmargin=*,itemsep=2pt,topsep=2pt]
  \item \textbf{Information layer.} We prove that the information about an aligned
    behavior $S$ lost by ablating a pathway $R_i$ is \emph{exactly} the conditional
    mutual information $\II(S;R_i\mid R_{-i})$, the part of $R_i$'s information that no
    other pathway carries (Theorem~\ref{thm:a1}). Fully redundant pathways cost zero.
    This is an exact identity, free of the LayerNorm artifact that contaminates the
    logit-level effect. A $k$-redundancy condition then guarantees resilience to any
    $k-1$ simultaneous ablations (Theorem~\ref{thm:a2}).
  \item \textbf{From information to behavior.} Through Fano's inequality, the
    information loss lower-bounds the unavoidable behavioral error of \emph{any}
    decoder reading the surviving pathways (Theorem~\ref{thm:b1}). Zero information
    loss forces zero behavioral error: redundancy implies provable behavioral
    robustness. Conversely, when behavior survives an ablation, the surviving pathways
    must have already contained all the information---compensation \emph{reveals}
    pre-existing redundancy rather than creating it (Proposition~\ref{prop:d1}).
  \item \textbf{Graph layer and phase transition.} Modeling computation as a layered
    directed acyclic graph, we show the \emph{targeted} critical ablation threshold
    equals the minimum vertex cut $R(G)$ (Theorem~\ref{thm:c1}), via Menger's theorem.
    Under random ablation, survival obeys the exact law $(1-q^w)^d$
    (Theorem~\ref{thm:c2}), with a percolation phase transition whose critical point
    $q_c=(1-2^{-1/d})^{1/w}\to 1$ and whose width shrinks as $1/w$
    (Theorem~\ref{thm:c3}). The targeted and random thresholds are therefore
    exponentially separated in the redundancy width.
\end{itemize}

The payoff is conceptual. A network carries \emph{two} robustness numbers: the
min-cut $R(G)$ (targeted fragility) and the percolation threshold $q_c$ (random
robustness), and Theorem~\ref{thm:c3} shows they diverge. The same aligned circuit is
then simultaneously robust to random or incidental damage and fragile to a targeted
attack on its min-cut. This reconciles the contradictory empirical reports and yields
an honest verdict on the premise that alignment intrinsically resists removal: the
premise holds in the random/benign regime but fails against a targeted adversary, who
need only degrade $R(G)$ pathway-units to misalign the model
(Corollary~\ref{cor:e1})---matching the easy targeted misalignment of
\cite{qi2023finetuning,qi2024shallow}.

\paragraph{Techniques.}
The proofs use only standard, uncontested tools: the chain rule and non-negativity of
(conditional) mutual information \cite{cover2006elements}, Fano's inequality
\cite{cover2006elements}, Menger's theorem in its vertex form \cite{diestel2017graph},
and exact percolation algebra on a series-parallel circuit. We deliberately base the
main bounds on Shannon mutual information rather than a specific partial-information-decomposition
operator, using the latter \cite{williams2010pid} only as an independent cross-check.

\paragraph{Organization.}
Section~\ref{sec:prelim} fixes notation, gives the definitions, and recalls the
prerequisite theorems. Section~\ref{sec:results} states and proves the information-layer
results (Section~\ref{sec:results}.1), the behavioral and structural corollaries
(Section~\ref{sec:results}.2), and the graph-layer results and phase transition
(Section~\ref{sec:results}.3), with worked examples and a computational-verification
subsection. Section~\ref{sec:discussion} discusses implications, reconciles the
literature, and lists open problems; Section~\ref{sec:conclusion} concludes.
